\section{Interpretation}

\begin{frame}{Ruling out simple alternatives}
  \small
  \begin{columns}[T,onlytextwidth]
    \begin{column}{0.55\textwidth}
      \begin{block}{Alternatives discussed in the paper}
        \begin{itemize}
          \item Incomplete or unrepresentative networks
          \item “Exceptional women” rather than a broad pattern
          \item Role monopolisation by women
          \item Over-reading a small set of centrality measures
          \item Online attention effects in male-dominated spaces
        \end{itemize}
      \end{block}
    \end{column}

    \begin{column}{0.45\textwidth}
      \begin{exampleblock}{Discussion points}
        \begin{itemize}
          \item ISIS groups can approach 50\% women, so rarity alone is weak.
          \item In PIRA, expert reputation lists did not align with top-BC women.
          \item No clear evidence ties BC to a single operational role.
        \end{itemize}
      \end{exampleblock}
    \end{column}
  \end{columns}

  \note{
    \textbf{Time: 2:15}\par
    Treat this as a credibility slide.\par
    Move quickly but explicitly through the alternatives the authors themselves consider.
  }
\end{frame}

\begin{frame}{Limitations and validity boundaries}
  \small
  \begin{block}{Data limitations}
    \begin{itemize}
      \item The authors cannot prove completeness; they argue the networks are representative enough to reveal real effects.
      \item External events shape network evolution and platform behaviour.
    \end{itemize}
  \end{block}

  \begin{block}{Measurement limitations}
    \begin{itemize}
      \item Online gender declarations may contain noise, though the paper argues misdeclaration incentives are limited.
      \item Higher-order network statistics need more data and are more sensitive to noise.
    \end{itemize}
  \end{block}

  {\footnotesize\color{PresTwoSecondary}\textbf{Inference discipline:}
  claims are about \alert{topological roles}, not intent or psychology.}

  \note{
    \textbf{Time: 2:15}\par
    Say clearly what the paper does not claim.\par
    Keep this slide as a boundary marker before discussing implications.
  }
\end{frame}

\begin{frame}{Implications: rethinking “importance” in networked systems}
  \small
  \begin{columns}[T,onlytextwidth]
    \begin{column}{0.55\textwidth}
      \begin{block}{Interpretive implication}
        Traditional approaches often focus on high-degree hubs; the paper argues that women's interconnectivity and \alert{collective measures like betweenness} also matter.
      \end{block}

      \begin{exampleblock}{Ethical constraint from the paper}
        The authors do \alert{not} recommend targeting women; they frame the implication as rethinking evaluation and engagement.
      \end{exampleblock}
    \end{column}

    \begin{column}{0.45\textwidth}
      \centering
      \begin{tikzpicture}[node distance=0.9cm]
        \node[draw=PresTwoMuted,rounded corners,fill=PresTwoSurface,align=center,inner sep=8pt] (hub) {Hub-centric\\\small degree};
        \node[draw=PresTwoMuted,rounded corners,fill=PresTwoSurface,align=center,inner sep=8pt,below=of hub] (broker) {Broker-centric\\\small betweenness};
        \draw[-Latex,thick] (hub) -- (broker);
      \end{tikzpicture}
    \end{column}
  \end{columns}

  \note{
    \textbf{Time: 1:30}\par
    Keep the implication conceptual, not tactical.\par
    This slide should sound like “how to evaluate importance,” not “how to exploit a network.”\par
    Then frame the model as a speculative add-on rather than core evidence.
  }
\end{frame}

\begin{frame}{A speculative mechanism: team ethic in a fission--fusion network}
  \small
  \begin{block}{Working assumptions}
    \centering
    \begin{tikzpicture}[node distance=0.45cm,scale=0.92,every node/.style={transform shape}]
      \node[draw=PresTwoMuted,rounded corners,fill=PresTwoSurface,align=center,inner sep=5pt,text width=0.23\textwidth] (a1) {High-BC actors are modeled as\\more team-oriented};
      \node[draw=PresTwoMuted,rounded corners,fill=PresTwoSurface,align=center,inner sep=5pt,text width=0.29\textwidth,right=of a1] (a2) {That team ethic spreads\\within clusters containing women};
      \node[draw=PresTwoMuted,rounded corners,fill=PresTwoSurface,align=center,inner sep=5pt,text width=0.18\textwidth,right=of a2] (a3) {The process runs atop a\\fission--fusion model};
      \draw[-Latex,thick] (a1)--(a2)--(a3);
    \end{tikzpicture}
  \end{block}

  \begin{columns}[T,onlytextwidth]
    \begin{column}{0.6\textwidth}
      \PaperFigureSlot[0.42\textheight]{2E.png}{
        Show both model-vs-data panels together.
      }
    \end{column}

    \begin{column}{0.4\textwidth}
      \begin{exampleblock}{Why include this at all?}
        PIRA observations are yearly, while the proposed spreading dynamics are assumed to happen faster. The paper presents this as an \alert{oversimplified, notional mechanism}, not as causal proof.
      \end{exampleblock}
    \end{column}
  \end{columns}

  \note{
    \textbf{Time: 2:30}\par
    Be explicit that the authors call this setup oversimplified.\par
    Explain fission--fusion only at a high level.\par
    Then move directly to the takeaways.
  }
\end{frame}

\begin{frame}{Conclusions and open questions}
  \begin{block}{Three take-home messages}
    \begin{enumerate}
      \item In both the online and offline cases, women show higher \alert{betweenness centrality} despite being fewer in number.
      \item Degree and betweenness can diverge: women can avoid visible “star” status while providing brokerage.
      \item Longevity patterns are consistent with women's embeddedness having system-level consequences.
    \end{enumerate}
  \end{block}

  \begin{exampleblock}{Open questions}
    How stable is this pattern across other extreme networks? Which pressures change centrality patterns? How should contributions be evaluated: locally or collectively?
  \end{exampleblock}

  \note{
    \textbf{Time: 0:45}\par
    Deliver these as compact sentences.\par
    This is the memory slide immediately before Q\&A.
  }
\end{frame}

\begin{frame}[plain]
  \vspace{1.1cm}
  \begin{center}
    {\fontsize{34}{38}\selectfont \textbf{Questions}}
  \end{center}

  \note{
    \textbf{Time: 7:00}\par
    If questions are short, use the backup slides for definitions, null model logic, or data scales.\par
    If questions are long, stop on time and offer to share the deck afterward.
  }
\end{frame}
